\subsection{RNA-count strata bound the interpretation of rank transport}
A frozen descriptive audit reuses the admitted target counts and source-only predictions, without new fitting. All 5,110 genes are ordered by RNA count and gene ID and split into four equal-count groups of 1,278, 1,278, 1,277 and 1,277. These are measurement-count strata, not biological replicates or a modeled noise estimate. Their RNA ranges are 98--6,899, 6,905--13,204, 13,210--26,317 and 26,318--2,752,618. Source-overlap counts are 1,135, 1,231, 1,231 and 1,158.

Within those groups, Spearman correlations against the raw FP/RNA log ratio are 0.186, 0.223, 0.283 and 0.179 for length/GC; 0.359, 0.462, 0.510 and 0.417 for codon frequencies; and 0.355, 0.412, 0.459 and 0.363 for the combined model. The codon-frequency model remains positive in every stratum, but its all-target correlation 0.526 is not a uniform within-stratum effect. Conditioning on a denominator count changes both the distribution and the question; it is not causal adjustment for RNA abundance.

Across all genes, codon-frequency predictions correlate 0.317 with RNA counts and 0.465 with footprint counts separately. Combined-model correlations are 0.260 and 0.402, while length/GC gives $-0.150$ and $-0.040$. Thus the ratio association coexists with associations to its components; it cannot be read as a direct intervention on translation efficiency. Library-size scaling, reused gene identities, lack of replicate uncertainty and the single WT pair remain limitations. No absolute protein yield, optimized synonymous sequence benefit or independent-patient/organism validation is inferred. Exact bins, all models and hashes are in \path{results/gse63789_count_strata_audit.json}. Source: \url{https://www.ncbi.nlm.nih.gov/geo/query/acc.cgi?acc=GSE63789}.
